\documentclass[a4paper,oneside,12pt]{extreport}

\usepackage{mmap}
\usepackage[T2A]{fontenc}
\usepackage[utf8]{inputenc}
\usepackage[english,russian]{babel}

% Текст отчёта следует печатать, соблюдая следующие размеры полей:
% левое — 30 мм, правое — 15 мм, верхнее и нижнее — 20 мм.
\usepackage[left=20mm, right=15mm, top=15mm, bottom=15mm]{geometry}

% \setlength{\parindent}{1.25cm} % Абзацный отступ

\usepackage{setspace}
%\onehalfspacing % Полуторный интервал

\frenchspacing % Равномерные пробелы
\usepackage{indentfirst} % Красная строка

\usepackage{microtype}
\sloppy

\usepackage{titlesec}
\titlespacing*{\chapter}{0pt}{-30pt}{8pt}
\titlespacing*{\section}{\parindent}{*4}{*4}
\titlespacing*{\subsection}{\parindent}{*4}{*4}
\titleformat{\chapter}{\LARGE\bfseries}{\thechapter}{20pt}{\LARGE\bfseries}
\titleformat{\section}{\Large\bfseries}{\thesection}{40pt}{\Large\bfseries}

\usepackage{graphicx}
\usepackage{caption}

\usepackage[unicode,pdftex]{hyperref}
\hypersetup{hidelinks}

%% title begin
\usepackage{wrapfig}

\makeatletter
	\def\vhrulefill#1{\leavevmode\leaders\hrule\@height#1\hfill \kern\z@}
\makeatother
%% title end

%% begin code
\usepackage{listings}
\usepackage{xcolor}

\lstset{
	basicstyle=\footnotesize\ttfamily,
	breakatwhitespace=true,
	breaklines=true,
	commentstyle=\color{gray},
	frame=single,
	keywordstyle=\color{blue},
	showstringspaces=false,
	stringstyle=\color{red},
	tabsize=8
}

\lstdefinestyle{lispinline}{
	frame=none,
	language=Lisp
}

\newcommand{\code}[1]{\texttt{#1}}
%% end code

%% begin theorem
\usepackage{amsthm}

\makeatletter
\newtheoremstyle{indented}
	{}% measure of space to leave above the theorem
	{}% measure of space to leave below the theorem
	{}% name of font to use in the body of the theorem
	{\parindent}% measure of space to indent
	{\bfseries}% name of head font
	{.}% punctuation between head and body
	{ }% space after theorem head; " " = normal interword space
	{}% header specification (empty for default)
\makeatother

\theoremstyle{indented}

\newtheorem{definition}{Определение}[section]
\newtheorem{example}{Пример}[section]
\newtheorem{theorem}{Теорема}[section]
\newtheorem{task}{Задание}

\makeatletter
\DeclareRobustCommand\bfseriesitshape{%
	\not@math@alphabet\itshapebfseries\relax
	\fontseries\bfdefault
	\fontshape\itdefault
	\selectfont
}
\makeatother

\DeclareTextFontCommand{\textbfit}{\bfseriesitshape}
\DeclareTextFontCommand{\define}{\bfseriesitshape}
%% end theorem

%% begin columns
\usepackage{etoolbox,refcount}
\usepackage{multicol}

\newcounter{countitems}
\newcounter{nextitemizecount}
\newcommand{\setupcountitems}{%
	\stepcounter{nextitemizecount}%
	\setcounter{countitems}{0}%
	\preto\item{\stepcounter{countitems}}%
}
\makeatletter
\newcommand{\computecountitems}{%
	\edef\@currentlabel{\number\c@countitems}%
	\label{countitems@\number\numexpr\value{nextitemizecount}-1\relax}%
}
\newcommand{\nextitemizecount}{%
	\getrefnumber{countitems@\number\c@nextitemizecount}%
}
\newcommand{\previtemizecount}{%
	\getrefnumber{countitems@\number\numexpr\value{nextitemizecount}-1\relax}%
}
\makeatother
\newenvironment{AutoMultiColItemize}{%
	\ifnumcomp{\nextitemizecount}{>}{3}{\begin{multicols}{2}}{}%
		\setupcountitems\begin{itemize}}%
		{\end{itemize}%
		\unskip\computecountitems\ifnumcomp{\previtemizecount}{>}{3}{\end{multicols}}{}}
\makeatother
\newenvironment{AutoMultiColEnumerate}{%
	\ifnumcomp{\nextitemizecount}{>}{3}{\begin{multicols}{2}}{}%
		\setupcountitems\begin{enumerate}}%
		{\end{enumerate}%
		\unskip\computecountitems\ifnumcomp{\previtemizecount}{>}{3}{\end{multicols}}{}}
%% end columns

\usepackage{array}
\usepackage{float}
\restylefloat{table}

\usepackage{amsmath}
\usepackage{amssymb}

\begin{document}

\begin{titlepage}
	{\large % 14pt instead of 12pt
	\onehalfspacing
	\centering

	\begin{wrapfigure}[7]{l}{0.14\linewidth}
		\vspace{3mm}
		\hspace{-10mm}
		\includegraphics[width=0.93\linewidth]{inc/img/bmstu-logo}
	\end{wrapfigure}
	{\singlespacing \footnotesize \bfseries Министерство науки и высшего образования Российской Федерации\\Федеральное государственное бюджетное образовательное учреждение\\высшего образования\\<<Московский государственный технический университет\\имени Н.~Э.~Баумана\\ (национальный исследовательский университет)>>\\(МГТУ им. Н.~Э.~Баумана)\\}

	\vspace{-2.2mm}
	\vhrulefill{0.9mm}\\
	\vspace{-7.5mm}
	\vhrulefill{0.2mm}\\
	\vspace{2mm}

	{\doublespacing \small \raggedright ФАКУЛЬТЕТ \hspace{25mm} «Информатика и системы управления»\\
	КАФЕДРА \hspace{5mm} «Программное обеспечение ЭВМ и информационные технологии»\\}

	\vspace{30mm}

	\textbf{ОТЧЁТ}\\
	По лабораторной работе №6\\
	По курсу: «Функциональное и логическое программирование»\\
	%Тема: «Тема работы»\\

	\vspace{60mm}

	\hspace{70mm} Студент:       \hfill Керимов~А.~Ш.\\
	\hspace{70mm} Группа:        \hfill ИУ7-64Б\\
	\hspace{70mm} Преподаватели: \hfill Толпинская~Н.~Б.,\\
	                             \hfill Строганов~Ю.~В.\\

	\vfill

	Москва\\
	\the\year\\}
\end{titlepage}

\setcounter{page}{2}


\begin{task}
	Пусть \code{(setf lst1 '(a b)) (setf lst2 '(c d))}.
	Каковы результаты вычисления следующих выражений?

	\begin{enumerate}
		\item \begin{lstlisting}[style=lispinline, gobble=24]
			(cons lst1 lst2) ; ((A B) C D)
		\end{lstlisting}

		\item \begin{lstlisting}[style=lispinline, gobble=24]
			(list lst1 lst2) ; ((A B) (C D))
		\end{lstlisting}

		\item \begin{lstlisting}[style=lispinline, gobble=24]
			(append lst1 lst2) ; (A B C D)
		\end{lstlisting}
	\end{enumerate}
\end{task}

\begin{task}
	Каковы результаты вычисления следующих выражений?

	\begin{AutoMultiColEnumerate}
		\item \begin{lstlisting}[style=lispinline, gobble=24]
			(reverse ()) ; NIL
		\end{lstlisting}

		\item \begin{lstlisting}[style=lispinline, gobble=24]
			(last ()) ; NIL
		\end{lstlisting}

		\item \begin{lstlisting}[style=lispinline, gobble=24]
			(reverse '(a)) ; (A)
		\end{lstlisting}

		\item \begin{lstlisting}[style=lispinline, gobble=24]
			(last '(a)) ; (A)
		\end{lstlisting}

		\item \begin{lstlisting}[style=lispinline, gobble=24]
			(reverse '((a b c))) ; ((A B C))
		\end{lstlisting}

		\item \begin{lstlisting}[style=lispinline, gobble=24]
			(last '((a b c))) ; ((A B C))
		\end{lstlisting}
	\end{AutoMultiColEnumerate}
\end{task}

\begin{task}
	Написать, по крайней мере, два варианта функции, которая возвращает последний элемент своего списка-аргумента.

	\begin{lstlisting}[language=Lisp, gobble=16]
		(defun my-last1 (x)
		  (car (reverse x)))

		(defun my-last2 (x)
		  (cond ((null (cdr x)) (car x))
		        (T (my_last2 (cdr x)))))

		(defun my-last3 (x)
		  (reduce (lambda (p n) n)))
	\end{lstlisting}
\end{task}

\begin{task}
	Написать, по крайней мере, два варианта функции, которая возвращает свой список-аргумент без последнего элемента.

	\begin{lstlisting}[language=Lisp, gobble=16]
		(defun lastout1 (x)
		  (reverse (cdr (reverse x))))

		(defun lastout2 (x)
		  (butlast x))
	\end{lstlisting}
\end{task}

\begin{task}
	Написать простой вариант игры в кости, в котором бросаются две правильные кости.
	Если сумма выпавших очков равна 7 или 11 — выигрыш, если выпало (1, 1) или (6, 6) — игрок получает право снова бросить кости, во всех остальных случаях ход переходит ко второму игроку, но запоминается сумма выпавших очков.
	Если второй игрок не выигрывает абсолютно, то выигрывает тот игрок, у которого больше очков.

	\begin{lstlisting}[language=Lisp, gobble=16]
		(defun throw-one () (+ (random 6) 1))

		(defun throw-two () (list (throw-one) (throw-one)))

		(defun next-try (dices)
		  (or (equal dices '(1 1))
		      (equal dices '(6 6))))

		(defun sum-dices (dices)
		  (+ (car dices) (cadr dices)))

		(defun is-win (dices)
		  (let ((rate (sum-dices dices)))
		       (or (= rate 7) (= rate 11))))

		(defun play-round ()
		  (let ((dices (throw-two)))
		       (cond
		         ((is-win dices) (list 0 dices))
		         ((next-try dices) (play-round))
		         (T (list (sum-dices dices) dices)))))

		(defun print-sum (d1 d2 msg)
		  (format nil "Player 1 - ~{~a~^ ~}~%Player 2 - ~{~a~^ ~}~%~a"
		    (cdr d1) (cdr d2) msg))

		(defun play ()
		  (let ((d1 (play-round))
		        (d2 (play-round)))
		       (cond
		         ((= (car d1) 0) (print-sum d1 d2 "Player 1 won!"))
		         ((= (car d1) 0) (print-sum d1 d2 "Player 2 won!"))
		         ((> (car d1) (car d2)) (print-sum d1 d2 "Player 1 won (sum)!"))
		         ((< (car d1) (car d2)) (print-sum d1 d2 "Player 2 won (sum)!"))
		         (T (print-sum d1 d2 "Draw!")))))
	\end{lstlisting}
\end{task}

\end{document}
